\chapter{序論}

\section{本章の概要}

本章では、研究の背景、問題提起、研究目的、研究課題、方法論の要点、本研究の貢献、および本論文の構成を述べる。

\section{研究背景}

Physical AI の発展により、AI は画像や文章を理解するだけでなく、物理世界で行動するロボットへ接続されつつある。特に、VLA は視覚、言語、行動を統合し、汎用的な操作能力をロボットへ与える方向として注目されている。RT-2 や OpenVLA に代表される研究は、短期的な操作を幅広い指示から実行する可能性を示している~\cite{rt2, openvla, rt1}.

しかし、実環境で必要となる作業は、単発の把持や移動だけではない。机上片付け、収納、配膳、掃除補助のような長期ロボットタスクでは、複数の工程を順に進める必要がある。その過程では、対象物の有無、遮蔽、所有者、危険性、ユーザ意図、作業履歴との不一致など、多様な不確実性が発生する。したがって、ロボットには「次に何をするか」だけでなく、「いま何が分かっていないのか」「その不確実性をどう扱うべきか」を判断する能力が求められる。

\section{問題提起}

既存研究では、不確実な場合に人間へ質問する研究、追加観測により見えない対象を確認する研究、失敗を検出する研究、安全上危険な動作を抑制する研究が個別に進められている~\cite{knowno, apvlm, aha}. しかし、長期タスク中では、同じ不確実性でも工程文脈によって適切な処理が変わる。

例えば、机上片付けにおいて「コップの中身が見えない」という不確実性は、コップを動かす直前であれば安全上の理由から stop または ask\_user が必要である。一方で、現在の工程がペンの片付けであり、コップ操作が後工程であれば defer が妥当である。また、「対象物が見えない」という状況も、探索中なら reobserve、まだ未実行なら act、配置済み履歴があるなら reobserve へ変化しうる。したがって、長期ロボットタスクでは、不確実性の有無だけでなく、その不確実性が現在工程においてどの処理分岐へ接続されるべきかを診断する必要がある。

\section{研究目的と対象範囲}

本研究の目的は、長期ロボットタスクにおいて、VLM を用いて工程文脈依存的な不確実性を診断し、現在の不確実性を act / reobserve / ask\_user / stop / defer のいずれへ接続すべきかを判定する方法を検討することである。

対象タスクは、机上片付けを想定した長期ロボットタスクとする。対象物はコップ、ペン、小物、布、箱、ケーブル、スポンジブロックなどを含み、各工程には「見つかっているか」「安全に持てるか」「確認が済んでいるか」などの工程条件が存在する。本研究では、大規模な world model や task and motion planning そのものを新規提案することは目的とせず、あくまで次工程へ進む前の工程条件不確実性を診断する前段モジュールに焦点を当てる。

\section{研究課題}

本研究では、以下の三点を主要課題として設定する。
\begin{enumerate}
  \item 長期ロボットタスク中の不確実性を、工程文脈と処理分岐の観点から整理する。
  \item VLM による不確実性診断手法を構成し、証拠充足性、不確実性原因、工程関連性、処理分岐を構造化して出力させる。
  \item 提案手法が実際の長期タスクにおいて有効かを、同一視覚状態・異なる工程文脈の対照評価と実機デモにより検証する。
\end{enumerate}

\section{方法論}

本研究では、現在画像、作業履歴、現在工程、次工程候補、工程条件リストを VLM へ入力する。VLM は各工程条件について、証拠が十分か、不確実性の原因は何か、現在工程において重要か、どの処理分岐が適切かを構造化形式で出力する。提案手法の特徴は、単に「次にどうするか」を直接出力させるのではなく、中間的な診断を経由して最終分岐を決定する点にある。

\section{本研究の貢献}

本研究の貢献は以下の三点である。
\begin{enumerate}
  \item 長期ロボットタスク中の不確実性を、工程文脈に応じた処理分岐として整理する。
  \item VLM を用いて、証拠充足性、不確実性原因、工程関連性、処理分岐を構造化して診断する手法を提案する。
  \item 同一視覚状態・異なる工程文脈の対照評価により、VLM が文脈依存的に act / reobserve / ask\_user / stop / defer を切り替えられるかを検証する。
\end{enumerate}

\section{本論文の構成}

本論文の構成は以下のとおりである。
第2章では関連研究を述べる。
第3章では提案手法を述べる。
第4章ではデータと実験設計を述べる。
第5章では実験1を中心とした分類性能評価を述べる。
第6章では実験2および実験3を通じた文脈依存性評価と実機デモを述べる。
第7章では考察を述べる。
第8章では結論と今後の課題を述べる。
